% Source: Wolf (2012), Chapter 8, upper-triangular power lemma

\begin{theorem}[Noncommutative upper-triangular word bound]
    \label{thm:wolf_upper_triangular_word_bound}
    \lean{Matrix.wolf_eq_105_seminorm}
    \leanok
    Let $D\ge 1$, let $\Lambda,N\in M_D(R)$ for a ring $R$, assume that $\Lambda$ is
    diagonal and that $N$ is strictly upper triangular, and let $\nu$ be an
    arbitrary submultiplicative ring seminorm.  Then, for every
    $n\in\mathbb N$,
    \begin{align}
      \nu\bigl((\Lambda+N)^n\bigr)
      \le \nu(\Lambda^n)
      +\sum_{k=1}^{\min\{n,D-1\}}
        \binom{n}{k}\nu(N)^k\nu(\Lambda)^{n-k}.
    \end{align}
    In particular, this applies to every submultiplicative norm in Wolf's
    upper-triangular power lemma
    \cite[Chapter~8, Equation~(8.105)]{Wolf2012Quantum}.
\end{theorem}

\begin{proof}
    \leanok
    Expand $(\Lambda+N)^n$ over all ordered words of length $n$ in the two
    letters $\Lambda$ and $N$; no commutativity is assumed.  A diagonal factor
    preserves the matrix index, whereas each strictly upper-triangular factor
    forces a strict increase.  Consequently every word containing at least
    $D$ copies of $N$ vanishes.  There are $\binom{n}{k}$ words containing
    exactly $k$ copies of $N$.  Subadditivity and submultiplicativity of $\nu$
    bound every such word by $\nu(N)^k\nu(\Lambda)^{n-k}$, giving the stated
    sum.  The word containing no copy of $N$ is kept exactly as
    $\nu(\Lambda^n)$, so no normalization condition on $\nu(1)$ is needed.
\end{proof}

\begin{theorem}[Upper-triangular power estimate]
    \label{thm:wolf_upper_triangular_power_estimate}
    \lean{Matrix.wolf_eq_103, Matrix.wolf_eq_103_refined}
    \leanok
    \uses{thm:wolf_upper_triangular_word_bound}
    Under the hypotheses of Theorem~\ref{thm:wolf_upper_triangular_word_bound},
    suppose in addition that $\nu(\Lambda)\le 1$.  If $D-1\le n$, then
    \begin{align}
      \nu\bigl((\Lambda+N)^n\bigr)
      \le \nu(\Lambda^n)
      +(D-1)n^{D-1}\nu(\Lambda)^{n-D+1}
        \max\{\nu(N),\nu(N)^{D-1}\}.
    \end{align}
    If $2(D-1)\le n$, the factor $(D-1)n^{D-1}$ may be replaced by
    $(D-1)\binom{n}{D-1}$.
\end{theorem}

\begin{proof}
    \leanok
    \uses{thm:wolf_upper_triangular_word_bound}
    Apply Theorem~\ref{thm:wolf_upper_triangular_word_bound}.  For
    $1\le k\le D-1$, bound $\binom{n}{k}$ by $n^{D-1}$ and bound
    $\nu(N)^k$ by $\max\{\nu(N),\nu(N)^{D-1}\}$.  Since
    $0\le\nu(\Lambda)\le1$ and $D-1\le n$,
    \begin{align}
      \nu(\Lambda)^{n-k}\le\nu(\Lambda)^{n-D+1}.
    \end{align}
    There are $D-1$ possible positive values of $k$, which gives the first
    constant.  If $2(D-1)\le n$, the binomial coefficients are increasing for
    $1\le k\le D-1$, so each is at most $\binom{n}{D-1}$.

% Wolf prints the first estimate for every $n\in\mathbb N$, although the
% exponent $n-D+1$ is negative when $n<D-1$ and is undefined at
% $\nu(\Lambda)=0$ without an additional convention.  The estimate above is
% therefore stated in the range $D-1\le n$, where the displayed
% natural-power derivation is meaningful.
\end{proof}

\begin{lemma}[Operator norm of matrix powers]
    \label{lem:l2_operator_norm_matrix_power}
    \lean{Matrix.l2_opNorm_pow_le_pow}
    \leanok
    For every $A\in M_D(\mathbb C)$ with $D\ge1$ and every $n\in\mathbb N$,
    \begin{align}
      \lVert A^n\rVert_\infty\le\lVert A\rVert_\infty^n.
    \end{align}
\end{lemma}

\begin{proof}
    \leanok
    For positive powers this is submultiplicativity.  For the zeroth power,
    the identity matrix has operator norm one because $D\ge1$.
\end{proof}

\begin{theorem}[Schur-form spectral-radius estimate]
    \label{thm:wolf_schur_form_spectral_radius_estimate}
    \lean{Matrix.wolf_eq_111_schur_form,
      Matrix.wolf_eq_111_schur_form_refined}
    \leanok
    \uses{thm:wolf_upper_triangular_power_estimate,
      lem:l2_operator_norm_matrix_power}
    Let $\Lambda,N\in M_D(\mathbb C)$, where $\Lambda$ is diagonal and $N$ is
    strictly upper triangular.  Suppose
    $\lVert\Lambda\rVert_\infty=\mu\le1$.  If $D-1\le n$, then
    \begin{align}
      \lVert(\Lambda+N)^n\rVert_\infty
      \le \mu^n
      +(D-1)n^{D-1}\mu^{n-D+1}
        \max\{\lVert N\rVert_\infty,\lVert N\rVert_\infty^{D-1}\}.
    \end{align}
    If $2(D-1)\le n$, the factor $(D-1)n^{D-1}$ may be replaced by
    $(D-1)\binom{n}{D-1}$.
\end{theorem}

\begin{proof}
    \leanok
    \uses{thm:wolf_upper_triangular_power_estimate,
      lem:l2_operator_norm_matrix_power}
    Apply Theorem~\ref{thm:wolf_upper_triangular_power_estimate} to the
    $\ell^2$ operator norm, use submultiplicativity to bound
    $\lVert\Lambda^n\rVert_\infty\le\lVert\Lambda\rVert_\infty^n$, and
    substitute the norm of the diagonal part.
\end{proof}
